\section{Minimizing enclosures and their area envelope}
\label{q:geometry:section}

The area in the numerical inequality is an infimum over complete
enclosing frontiers. A metric perturbation can change which frontier
attains that infimum. We first describe the whole family of minimizers
and then differentiate its least area. This argument concerns only the
metric and the enclosing sets; no constraint equation or equality
hypothesis enters it.

Let $(X,g)$ be a connected orientable smooth $n$-manifold, $n\ge3$,
complete with its nonempty compact smooth boundary $B$ included.
The boundary may be disconnected. Assume that outside a compact set
$X$ is one Euclidean coordinate end and
\begin{equation}\label{q:geometry:AF}
 g-\delta=O_2(r^{-q}),\qquad q>0.
\end{equation}
A full cut is the entire compact smooth embedded two-sided intrinsic
boundary of a connected closed outer domain in $X$, whose interior
lies in $\operatorname{int}X$ and which contains the distant end.
Every boundary component and every portion coincident with $B$ is
counted. Write $a(g)$ for the infimum of their full areas.

Attach a compact smooth filling $O$ behind $B$ and extend the metric
smoothly to the resulting boundaryless manifold $\widehat X$. Such a
filling exists here: reverse a compact truncation of $X$ and cap its
spherical end boundary with a ball. The remaining boundary is the
oppositely oriented $B$. The filling and collar extension serve only
to measure perimeter on $B$; the metric farther inside $O$ has no role.
Let $\mathcal A$ be the bounded finite-perimeter sets $E\subset\widehat X$
containing $O$ almost everywhere, identified modulo null sets. Put
\begin{equation}\label{q:geometry:full-perimeter}
 P_g(E)=|D\mathbf1_E|_g(\widehat X),\qquad
 T_E=\operatorname{supp}|D\mathbf1_E|_g,
 \qquad \mathcal T(g)=\operatorname*{argmin}_{E\in\mathcal A}P_g(E).
\end{equation}
Thus the perimeter is measured in the filled manifold, including
contact with $B$. A smooth full cut bounds a filled set of precisely
the same perimeter; in particular $P_g(O)=\operatorname{Area}_g(B)$.
We use the basic BV compactness, Gauss--Green and submodularity facts
in \cite{AmbrosioFuscoPallara2000,Maggi2012}.

\subsection{Existence, contact regularity, and strict barriers}

\begin{proposition}[The full minimizing enclosure]
\label{q:geometry:enclosure}
The minimum in Equation~\eqref{q:geometry:full-perimeter} exists,
equals $a(g)>0$, and is independent of the filling. All its minimizing
sets are contained in one compact region. Their multiplicity-one
frontiers are smooth minimal hypersurfaces away from $B$ and a
singular set of Hausdorff dimension at most $n-8$; the singular set is
empty when $n\le7$. Every contact point with $B$ is regular, with a
$C^{1,\alpha}\cap W^{2,p}$ graph for every finite $p$ and
$0<\alpha<1$. For $n=3,4$ the whole frontier is $C^{1,1}$, including
contact, and its graph functions are $W^{2,\infty}$.

In dimensions $3$ and $4$, each frontier bounds a connected exterior
and is the $C^1$ and area limit of smooth full cuts contained in
$\operatorname{int}X$. If $H_B<0$ for the normal into $X$, every
minimizer contains one fixed exterior collar of $B$. In dimensions
$3$ and $4$ its frontier is then a smooth minimal full cut; its closed
exterior is connected, one-ended, complete with its boundary included,
and outer area-minimizing with all boundary components counted.
In every dimension, whenever a minimizing frontier is smooth, its
closed exterior has these same connectedness, completeness, one-end
and outer area-minimization properties.
\end{proposition}

\begin{proof}
\emph{Confinement and existence.}
For sufficiently large $R_0$ the end spheres have
\[
 H_{\{r=R\}}=\operatorname{div}_g
       (\nabla r/|\nabla r|_g)
       =\frac{n-1}{R}+O(R^{-1-q})>0.
\]
Set $Y=\nabla r/|\nabla r|_g$. For each competitor $E$, and almost
every $R>R_0$, the BV trace formula and Gauss--Green on
$E\cap\{r>R\}$ give
\begin{equation}\label{q:geometry:clipping}
 P_g(E\setminus\{r>R\})-P_g(E)
 \le-\int_{E\cap\{r>R\}}\operatorname{div}_gY\,dV_g\le0.
\end{equation}
Indeed the new slice area is the flux through the inner end of this
region; the original exterior-boundary flux is at most its perimeter.
At a good slicing radius the two traces agree, so these are exactly
the terms in the perimeter difference. The inequality is strict if
the removed set has positive volume.

Fix $R_1>R_0$. For each competitor choose a good radius in
$(R_0,R_1)$ and clip there. This confines a minimizing sequence to
the same compact region. BV compactness and lower semicontinuity
produce a global minimizer; containment of $O$ and emptiness in the
fixed tail are closed constraints. Applying the strict inequality to
any minimizer, at a good radius below a fixed $R_2$ with
$R_0<R_2<R_1$, puts every minimizer inside the smaller truncation.
There is no need for a radius that is simultaneously a good slice
for every competitor.

\emph{Recovering the full smooth-cut infimum.}
A full cut supplies a member of $\mathcal A$, so
$\min P_g\le a(g)$. In the other direction choose a smooth collar
field pointing strictly out of $O$ on $B$ and let $\Phi_\varepsilon$
be its flow. For small positive $\varepsilon$, the set
$\Phi_\varepsilon(E)$ contains a neighborhood of $O$, and its
perimeter tends to $P_g(E)$. Approximate its indicator in finitely
many charts by convolution and a partition of unity, keeping it
identically one on a smaller neighborhood of $O$ and zero outside a
fixed compact region. The BV approximation and coarea formula give
smooth level sets containing that neighborhood, with perimeters
tending to $P_g(\Phi_\varepsilon(E))$.
Fill every bounded complementary component of such a level set and
retain the complementary component containing the end. Its closure
in $X$ is a connected smooth outer domain; filling deletes boundary
components and creates none. Its full-cut area is no larger than the
approximating perimeter. Letting the approximation error, then
$\varepsilon$, tend to zero proves $a(g)\le P_g(E)$. This argument
also works when $E$ has singular frontier.

\emph{Local regularity.}
Choose a smooth field $Z$ supported near $B$, with $|Z|_g\le1$ and
$Z=\nu_O$ on $B$. For any local replacement $F$ of a minimizer $E$,
submodularity, constrained minimality, and Gauss--Green imply
\begin{align*}
 P_g(E)&\le P_g(F\cup O)
       \le P_g(F)+P_g(O)-P_g(F\cap O),\\
 P_g(O)-P_g(F\cap O)&\le
       \int_{O\setminus F}\operatorname{div}_gZ\,dV_g
       \le C\operatorname{Vol}_g(E\triangle F).
\end{align*}
The unchanged terms outside the replacement region cancel. Thus, if
$E\triangle F\Subset U$,
\begin{equation}\label{q:geometry:almost-minimality}
 P_g(E;U)\le P_g(F;U)+C\operatorname{Vol}_g(E\triangle F).
\end{equation}
Ball insertion and deletion, followed by the relative isoperimetric
inequality, give an upper perimeter bound $Cr^{n-1}$ and positive
lower density bounds for both phases at every frontier point
\cite{Maggi2012}. The
volume error in Equation~\eqref{q:geometry:almost-minimality} scales
as $r^n$. The local almost-minimizing-boundary theorem gives
$C^{1,\alpha}$ regularity at reduced-boundary points and minimizing
boundary cones as tangent sets. Dimension reduction gives singular
dimension at most $(n-1)-7=n-8$; see
\cite[Theorem~2.27]{BiZhu2026} and the codimension-one theory in
\cite[Chapter~7]{Simon2018}. Off the obstacle the minimizing property
has zero error, so the regular frontier is smooth and minimal.

At a contact point a tangent filled set contains the tangent
half-space of $O$. Its stationary boundary cone lies on one side of
the bounding plane. Such a cone is planar: the nonnegative height
function on its spherical link satisfies weakly
$\Delta_{\rm link}z=-(n-2)z$. Testing with the constant function
makes its height integral zero, hence its support lies in the plane.
The weak stationary identity is on the whole link and does not
discard its singular set. The two density bounds and the
multiplicity-one boundary property select a half-space rather than
the full or empty set. Density regularity therefore makes the contact
point a graph point.

Equation~\eqref{q:geometry:almost-minimality}, applied to flows with
both signs, bounds the graph's weak mean curvature. Its scalar
minimal-surface operator has smooth coefficients and is uniformly
elliptic on the bounded gradient range. Difference quotients yield
$W^{2,2}$ locally. In nondivergence form the leading coefficients
are $C^{0,\alpha}$ and the right side is bounded; localized linear
$W^{2,p}$ estimates then give every finite $p$. The Sobolev embedding
allows every $\alpha<1$. In dimensions $3$ and $4$ we also use the
stronger smooth-obstacle theorem: for a $C^2$ obstacle, zero bulk
term, and ambient dimension below eight, its entire minimizing
frontier is $C^{1,1}$
\cite[Regularity Theorem~1.3(iii), pp.~368--369]{HuiskenIlmanen2001}.
This gives the stated $W^{2,\infty}$ regularity.

The frontier is nonempty. Otherwise its BV indicator is constant
on the connected filled manifold, contrary to containment of the
open obstacle and absence from an end tail. Its perimeter is positive;
this follows from the relative isoperimetric inequality, or from a
regular frontier chart.

\emph{Connected exterior and smooth approximation in low dimensions.}
Take the closure of the measure-theoretic interior as representative.
When $n=3,4$ its frontier is compact embedded $C^{1,1}$, with filled
and exterior sets on its two local sides. Any complementary component
other than the end component is relatively compact. Its boundary is
a nonempty union of frontier components. Filling it removes positive
full perimeter, contradicting minimality. The exterior is connected. The same argument applies in every
dimension whenever the frontier is smooth. It only uses the local
two-sided hypersurface charts and compactness.

Choose a smooth field $V$ uniformly transverse to the continuous
outward normal of the compact frontier $T$. Its flow gives a $C^1$
collar $\Psi:T\times(-\delta,\delta)\to\widehat X$ with filled side
$t\le0$. For small $\varepsilon>0$, the flowed set is
$E\cup\Psi(T\times(0,\varepsilon])$; every crossing of the frontier
is outward. It contains $O$ in its interior, and the displaced
frontier lies at positive distance from $B$. Smooth a $C^1$ defining
function in a smaller collar, retaining its positive derivative
along $V$. The resulting smooth zero set is a small graph on the
collar fibers. Extend the graph displacement to an isotopy supported
in that collar; it fixes $O$ and the distant end. Enclosure and
connectedness of the exterior are therefore preserved. The resulting
full cuts converge in $C^1$ and area to $T$.

\emph{Detachment from a strict inner barrier.}
If $H_B<0$ for the normal into $X$, a fixed distance collar
$0\le t\le\varepsilon$ has $Y=\partial_t$ and
$\operatorname{div}_gY\le-\beta<0$. For a good level $s$, put
$O_s=O\cup\{0<t<s\}$ and $W_s=O_s\setminus E$.
Gauss--Green on $W_s$, including the flux on $t=0$, gives
\begin{equation}\label{q:geometry:detachment}
 P_g(E\cup O_s)-P_g(E)
 \le\int_{W_s}\operatorname{div}_gY\,dV_g
 \le-\beta\operatorname{Vol}_g(W_s).
\end{equation}
Minimality makes $W_s$ null. Choosing, for each minimizer, a good
$s\in(\varepsilon/2,\varepsilon)$ proves that the same smaller
collar lies in its filled interior. Thus all contact is removed.

For $n=3,4$ the detached frontier is now smooth and minimal. More
generally, in any dimension where the frontier under consideration
is smooth, its connected closed exterior $D$ contains just the original end. It is complete
in its intrinsic length metric: an intrinsic Cauchy sequence is
Cauchy in $X$, has a limit in the closed set $D$, and smooth interior
or boundary half-ball charts give intrinsic convergence to that
limit. Finally, every full enclosing cut of $D$ gives a filled
competitor containing $E$ and hence $O$. Its full area is at least
$P_g(E)$. This proves outer area-minimization with every component
and all coincidence counted.
\end{proof}

\subsection{Continuity of planes before differentiation}

For $E\in\mathcal T(g)$ define the position-and-plane measure $V_E$
on the bundle of unoriented $(n-1)$-planes by
\begin{equation}\label{q:geometry:plane-measure}
 \int\Phi\,dV_E=
 \int_{\partial^*E}\Phi(x,T_x\partial^*E)\,dA_g.
\end{equation}
Singular points have zero perimeter measure. Give $\mathcal T(g)$
the topology of $L^1$ convergence of indicators and weak convergence
of these measures. The next result explains why its second requirement
follows from the first, and why it is needed for the metric derivative.

\begin{proposition}[Compactness and the derivative of least area]
\label{q:geometry:envelope}
The space $\mathcal T(g)$ is compact and metrizable. For every smooth
symmetric two-tensor $h$ the function
\[
 E\longmapsto a'_E(h):=
 \frac12\int_{\partial^*E}\operatorname{tr}_{T\partial^*E}h\,dA_g
\]
is continuous on it.
Let $g_t$, $|t|<t_0$, be a smooth path of smooth metrics with $g_0=g$,
extended smoothly through the fixed boundary. Assume uniform
comparison with $g$, and bounded first and second parameter derivatives
relative to $g$. One sufficient end hypothesis is
\begin{equation}\label{q:geometry:uniform-AF}
 g_t-\delta=O_2(r^{-q_*})\quad\hbox{uniformly for }|t|<t_0,
 \qquad q_*>0.
\end{equation}
More generally, it suffices that the coordinate spheres have positive
outward mean curvature beyond one common radius. The result also
holds if the end metrics are uniformly comparable to the Euclidean
metric with uniform scaled $C^1$ bounds; they need not have a common
mean-convex tail in this last alternative. Put $h=\partial_tg_t|_0$. Then
\begin{equation}\label{q:geometry:envelope-derivative}
 \left.\frac{d}{dt}\right|_{0+}a(g_t)
       =\min_{E\in\mathcal T(g)}a'_E(h).
\end{equation}
For any $t_j\to0$ and any $E_j\in\mathcal T(g_{t_j})$, a subsequence
converges in $L^1$ to $E\in\mathcal T(g)$, with convergence of
perimeters and of the plane measures computed in the fixed metric $g$.
\end{proposition}

\begin{proof}
The uniform asymptotic hypothesis gives one mean-convex tail, so
Equation~\eqref{q:geometry:clipping} confines all minimizers uniformly.
The same proof applies under the common-sphere hypothesis. For the
last alternative, existence can be proved without an outer barrier.
Minimize inside successively larger closed coordinate truncations.
Each constrained minimum exists by the direct method and is at most
$P_{g_t}(O)$. Local BV compactness yields a limiting set containing
$O$, of finite total perimeter. The minimizing values decrease to the
infimum over bounded competitors. The limit is locally minimizing
away from $O$, and constrained-minimizing at $O$: to compare it with
a compact replacement, glue that replacement into the truncated
minimizers across a thin annulus. Coarea chooses a slicing boundary
on which the $L^1$ trace discrepancy tends to zero, so the added
perimeter tends to zero. The minimizing inequality then passes to
the limit. This is the local minimizing-boundary compactness argument
in the BV framework of \cite{AmbrosioFuscoPallara2000,Maggi2012}.

The limit has finite volume, so its exterior phase is the empty one
once its frontier is confined. To see the finite-volume assertion
before using confinement, apply the Euclidean isoperimetric inequality
to the filled part on the coordinate end with a fixed inner slicing
sphere, capping that slice inside the coordinate ball. Its Euclidean
perimeter is bounded by the original perimeter times the uniform
metric-comparison constant, plus the fixed sphere area. The resulting
volume bound is uniform in the outer truncation; local convergence
passes it to the limit.

Testing on $O$ gives a uniform perimeter bound. If a minimizing frontier passes through an
end point of radius $r$, a coordinate ball of radius $cr$ about that
point misses the obstacle. The rescaled metrics are uniformly elliptic
with uniformly bounded first derivatives, and the frontier is locally
perimeter-minimizing in that ball. The interior perimeter density
estimate gives area at least $c_1r^{n-1}$, with $c_1$ uniform. This
contradicts the fixed upper bound for large $r$. Thus the limiting
frontier lies in a fixed compact region. The connected end tail is
therefore almost everywhere full or empty; finite volume makes it
empty. The limiting set is an admissible bounded competitor. Lower
semicontinuity bounds its perimeter by the infimum approached by the
truncated minima, so it is a global minimizer. The same density argument
confines every global minimizer, uniformly in $t$. Thus their bounded
filled sides again lie in one compact set.

BV compactness now applies on that compact region. For fixed-metric
minimizers $E_j\to E$ in $L^1$, lower semicontinuity gives
\[
 a(g)\le P_g(E)\le\liminf_jP_g(E_j)=a(g).
\]
Thus the limit is a minimizer and the perimeters converge. Set
$\sigma_j=|D\mathbf1_{E_j}|_g$ and $\sigma=|D\mathbf1_E|_g$.
Weighted lower semicontinuity makes every weak limit of $\sigma_j$
dominate $\sigma$. Their total masses agree, so
$\sigma_j\rightharpoonup\sigma$. The normal vector measures
$\nu_j\sigma_j$ converge weakly to $\nu\sigma$ by Gauss--Green.

To pass from normals to tangent planes, cover the compact support by
finitely many charts with nonnegative smooth partition weights
$\rho_\alpha$. Choose a smooth matrix $A_\alpha(x)$ in each chart
such that $|A_\alpha(x)v|_\delta=|v|_g$. The Euclidean vector measures
\[
 \lambda_{\alpha,j}=\rho_\alpha A_\alpha\nu_j\sigma_j
 \quad\hbox{have}\quad
 |\lambda_{\alpha,j}|_\delta=\rho_\alpha\sigma_j.
\]
They converge weakly and their total variations have convergent
masses. The localized Reshetnyak continuity theorem
\cite{Reshetnyak1968}, in the polar-direction formulation of
\cite[Theorem~1.3]{Spector2011}, therefore applies to
\[
 (x,z)\longmapsto
 \Phi\bigl(x,(A_\alpha(x)^{-1}z)^{\perp_g}\bigr),
 \qquad |z|_\delta=1.
\]
Insert a continuous compact cutoff equal to one on the localized
support, so the integrand is bounded on the chart, and sum over the
partition. This proves weak convergence of $V_{E_j}$, including
contact mass: a chart crossing $B$ is an interior chart of
$\widehat X$. It is convergence of plane integrals, not yet pointwise
convergence of tangent planes. Compact supports and bounded masses
make the indicated measure topology metrizable. The subsequence
argument proves compactness, and the choice
$\Phi(x,\Pi)=\frac12\operatorname{tr}_\Pi h$ proves continuity.

For varying metrics, uniform convergence on the confining region
provides $\varepsilon_t\to0$ such that
\[
 (1-\varepsilon_t)P_g(F)\le P_{g_t}(F)
              \le(1+\varepsilon_t)P_g(F)
\]
for every confined competitor. Testing a fixed $g$-minimizer and
using $P_g(E_t)\ge a(g)$ shows
$a(g_t)\to a(g)$ and $P_g(E_t)\to a(g)$. Hence every BV limit is a
$g$-minimizer and the preceding localized argument gives its plane
measure convergence.

On the compact Grassmann bundle the area density has the uniform
Taylor expansion
\begin{equation}\label{q:geometry:area-taylor}
 P_{g_t}(F)=P_g(F)+t\,a'_F(h)+O(t^2)P_g(F).
\end{equation}
Testing every fixed $E\in\mathcal T(g)$ gives
$\limsup_{t\downarrow0}(a(g_t)-a(g))/t\le\min_Ea'_E(h)$.
For the other bound choose $E_t\in\mathcal T(g_t)$. Its $g$-perimeter
is uniformly bounded and at least $a(g)$, so
\[
 \frac{a(g_t)-a(g)}t\ge a'_{E_t}(h)-O(t).
\]
Every sequence tending to zero has a subsequence whose plane measures
converge to those of a $g$-minimizer. The right side then tends to a
value at least $\min_Ea'_E(h)$. This proves both the derivative and
its formula, with no uniqueness or smooth choice of a minimizing cut.
\end{proof}

\subsection{Tied sheets and uniform graphical separation}

The averaged area source later requires one plane at a point shared
by different minimizing frontiers. It also requires quantitative
control when a nearby frontier approaches a given regular sheet.
These are local consequences of minimizing geometry, distinct from
the integral continuity just proved.

\begin{proposition}[No crossing and uniform regular sheets]
\label{q:geometry:sheets}
If a minimizing frontier meets the free regular part of another,
the two frontiers agree locally at that point. In particular the
plane field on the union of their free regular parts is well-defined
and continuous in its relative topology. Around every such point
there are nested graph cylinders, disjoint from $B$, in which every
minimizing frontier meeting the smaller cylinder is a single smooth
graph over the same base. The graph bounds are uniform on smaller
base domains. Two such graphs either agree or are disjoint and can
be ordered. For $V'\Subset V$ and fixed $y_0\in V$, their ordered
heights obey
\begin{equation}\label{q:geometry:graph-separation}
 \|h_1-h_2\|_{C^j(V')}\le C_j|h_1(y_0)-h_2(y_0)|,
 \qquad j=0,1,2,\ldots,
\end{equation}
with constants uniform in the two minimizers.
\end{proposition}

\begin{proof}
For minimizing sets $E,F$, submodularity gives
\[
 P_g(E\cup F)+P_g(E\cap F)\le P_g(E)+P_g(F)=2a(g).
\]
Each term on the left is at least $a(g)$, so the union and
intersection are themselves minimizers and the submodularity deficit
is zero as a measure. Let $x\in T_F$ lie on a free regular sheet of
$T_E$. Blow up both sets at $x$ along a common sequence, writing
$H$ for the tangent half-space of $E$ and $C$ for a tangent set of
$F$. Their unions and intersections converge in $L^1$ to $C\cup H$
and $C\cap H$. Minimizing-boundary compactness makes these minimizing
cones, with empty and full sets allowed. The half-space argument in
Proposition~\ref{q:geometry:enclosure} gives
\[
 C\cup H\in\{H,\mathbb R^n\},\qquad
 C\cap H\in\{\varnothing,H\}.
\]
Thus $C$ is one of $H,H^c,\varnothing,\mathbb R^n$; the density
bounds for $F$ rule out the latter two.

If $C=H$, both union and intersection have density-one planar tangent
and are regular near $x$. Their minimal graphs are ordered against
the reference sheet and touch it at $x$. The strong comparison
principle makes them agree with it locally, hence so does $T_F$.
If $C=H^c$, the union has full tangent and the intersection empty
tangent. The uniform two-phase density bounds at frontier points
put $x$ outside the supports of both boundaries. In a neighborhood
the union is full and the intersection empty, so $F$ is the complement
of $E$ there. Their common open sheet contributes twice its positive
area to $P_g(E)+P_g(F)$ and zero to the union and intersection.
That contradicts the zero submodularity deficit. This excludes the
opposite coorientation and proves local agreement.

We next justify the uniform graphical assertion, including the
entire support. Suppose $x_j\in T_{E_j}$ tends to a free regular
point $x$ of a fixed minimizing frontier. Compactness from
Proposition~\ref{q:geometry:envelope} gives a minimizing limit
$E_\infty$. Uniform density bounds put $x\in T_{E_\infty}$, and
the preceding argument identifies its sheet near $x$ with the fixed
regular sheet. In a ball about $x$ missing $B$, all these boundaries
are locally area-minimizing. Their integral varifolds are stationary
in the smooth ambient metric, with multiplicity one and no boundary.
The limit is smooth with density one.

Choose a sufficiently small continuity radius for its area measure.
The limiting mass ratio is close to one and its height and tilt
excesses relative to $T_xT_{E_\infty}$ are small. Plane-measure
convergence transfers these properties to $E_j$. Moving centers to
$x_j$ changes the required radii by a quantity tending to zero.
To apply the small-excess theorem in a Riemannian chart, use its
smooth-ambient form, or smoothly isometrically embed a neighborhood:
stationarity in the ambient manifold gives Euclidean generalized
mean curvature bounded by the ambient second fundamental form.
After rescaling by the small radius, the required $L^p$ curvature
norm is small for any fixed $p>n-1$.

Allard's graphical theorem therefore represents the entire support
in a smaller ball by a single $C^{1,\alpha}$ graph with uniform bounds
\cite{Allard1972,Maggi2012}; the exact small-mass and $L^p$ mean-curvature
formulation is \cite[Chapter~5, Section~5, Theorem~5.2]{Simon2014}.
It leaves no second small sheet or spike outside a selected graphical
component. Uniform graph compactness and the identified varifold
limit give $C^1$ convergence; interior estimates for the smooth
minimal-graph equation improve this to smooth convergence on smaller
domains. A sequence contradicting uniform cylinders or uniform graph
bounds would have precisely such a convergent subsequence, which is
impossible. The same argument for $x_j\to x$ proves continuity of
the common plane field.

Two distinct local graphs do not meet by the first part of the proof,
so their heights can be ordered. Their nonnegative difference solves
a linear uniformly elliptic equation, obtained by integrating the
derivative of the scalar minimal-graph operator between the two
graphs. The preceding interior bounds give uniform coefficient bounds
of every order, including the lower-order terms. The Harnack
inequality for a nonnegative solution bounds its size on smaller
domains by its value at $y_0$; interior derivative estimates then give
Equation~\eqref{q:geometry:graph-separation}. If that value is zero,
the same principle makes the two graphs coincide.
\end{proof}
